\section{Receding-horizon control evaluation}

\subsection{Frozen regulation protocol}

The closed-loop study was specified before execution and reuses the frozen
PGNN-style and \method{} checkpoints. It includes all five systems, eight
seeds, and three plant/observation conditions: clean in-distribution (ID),
corrupted ID, and corrupted out-of-distribution parameters. Four independently
generated plant episodes populate each seed--system--condition cell. Initial
states, plant parameters, observation draws, masks, and planner draws are
paired across four prediction interfaces: nominal propagation from the last
imputed state, filter plus nominal propagation, PGNN-style, and \method. Clean
and corrupted ID evaluations share the plant draw.

Seven zero-input transitions establish the eight-state context. The controller
then regulates for 20 steps. Deterministic-mean cross-entropy model predictive
control uses a horizon of 8, 64 candidates, two iterations, and eight elites.
Action sequences start with standard deviation 0.6 of the action limit, with a
0.05 floor. The stage objective is mean squared state error in training-
standard-deviation units plus 0.01 times squared normalized action; terminal
state error receives an additional factor of two. The pendulum target is
$(\sin\theta,\cos\theta,\omega)=(0,1,0)$ and the other targets are zero.
Corrupted observations use Gaussian noise with scalar standard deviation 8\%
of the mean training-state standard deviation and 25\% independent channel
missingness.

The primary response is the realized mean normalized state loss over 20
controlled transitions. Episodes, systems, and conditions receive equal weight
within each seed. The reported \method{}--PGNN-style contrast is the mean of
eight seed-level relative improvements, with a 5,000-replicate paired seed
bootstrap interval and an exact one-sided sign-flip test against zero.
Controller settings were not changed after observing outcomes.

\subsection{Regulation outcomes}

Table~\ref{tab:closed-loop} reports \ClosedLoopEpisodes{} paired episodes.
CFR-WM obtains state loss \ClosedLoopCFRLoss{} versus
\ClosedLoopPGNNLoss{} for PGNN-style. The seed-balanced relative improvement
is \ClosedLoopImprovement\% (95\% paired bootstrap interval
[\ClosedLoopCILow\%, \ClosedLoopCIHigh\%]); the exact one-sided sign-flip
test against zero gives $p=\ClosedLoopPValue$. Mean improvements favor CFR-WM
in all 15 system--condition cells, ranging from 0.006\% for the damped-
oscillator out-of-distribution corrupted cell to 2.8\% for the cart--pole
corrupted-ID cell (Fig.~\ref{fig:closed-loop}).

\begin{table*}[t]
\centering
\caption{Closed-loop regulation with a common cross-entropy model predictive
controller. State loss and action effort are mean $\pm$ SD across eight
seed-level aggregates. Condition columns report state loss.}
\label{tab:closed-loop}
\resizebox{\textwidth}{!}{%
\input{table_closed_loop}}
\end{table*}

\begin{figure*}[t]
\centering
\includegraphics[width=0.92\textwidth]{figures/FigClosedLoop.pdf}
\caption{Frozen closed-loop study. (a) Seed-balanced mean normalized state
loss; error bars are one SD across eight seeds. (b) Mean CFR-WM improvement
over PGNN-style within each system--condition cell. Positive values favor
CFR-WM.}
\label{fig:closed-loop}
\end{figure*}

The separation between learned predictors is small. Filtering followed by
nominal dynamics obtains state loss \ClosedLoopFilterNominalLoss{}, slightly
below CFR-WM's \ClosedLoopCFRLoss{}, and uses less action effort (0.398 versus
0.404). The nominal-last interface reaches \ClosedLoopNominalLoss{}. The
closed-loop study therefore supports the value of context correction, but does
not show an additional benefit from the learned residual.
